\section{Proof of the Institutional Query-Complexity Theorem}
\label{sec:r19-complexity-proof}
\begin{proof}[Proof of Theorem~\ref{paper-thm:r19-query-complexity}]
For any $w\ne0$, the exact query $t=-hw$, $h>0$, returns $-h\min_{z\in Z}w\cdot z$. Thus $m$ queries on the stated rays identify all $m$ robust returns.

For necessity, run a deterministic adaptive acquisition rule against the economy $Z_r=\bar z+r\mathbb B^k$, with any fixed $0<r<R$. A nonzero query $t$ returns $t\cdot\bar z+r\|t\|$. If the procedure stops before querying $m$ distinct rays, some target direction $u=-w^j/\|w^j\|$ is not equal to any normalized queried direction. Finiteness then gives $a(u)<1$. Choose a sufficiently small $\delta>0$ so that $\delta\leq R-r$ and $(r+\delta)a(u)\leq r$. Define
\begin{equation}
 Z_{r,u,\delta}=\operatorname{co}\bigl(Z_r\cup\{\bar z+(r+\delta)u\}\bigr).
 \label{eq:r19-spike-set}
\end{equation}
The set is compact, convex, and contained in $\bar z+R\mathbb B^k\subset K$. For a normalized queried direction $q_i$, the added point has excess support $(r+\delta)q_i\cdot u\leq r$, and therefore does not alter the query answer. The zero query is also unchanged. At each step identical earlier answers lead the deterministic rule to the same next query, so induction shows that the complete adaptive transcript is unchanged. The robust return in target direction $w^j$ is nevertheless lower by $\delta\|w^j\|$. Exact identification of all robust returns is impossible. This proves the first part, including its equivalent all-safe-level interpretation.

The same construction proves the second part for every positive $\delta$ satisfying the displayed inequalities, not only for a small one. With $w=-u$, the robust returns are
\[
 b=-u\cdot\bar z-r\quad\text{and}\quad b-\delta.
\]
Set the safe offer's return to $b-\delta/2$. After the common transcript, let $p$ be the probability of choosing the response-sensitive offer. In the first economy regret is $(1-p)\delta/2$; in the second it is $p\delta/2$. Minimizing their maximum over $p\in[0,1]$ yields $p=1/2$ and regret $\delta/4$. This is an indistinguishability argument about information, not a bound on a numerical approximation residual.

For the third part, suppose the caps of radius $\alpha$ do not cover the sphere. There is a unit $u$ for which $a(u)<\cos\alpha=r/(r+4\varepsilon)$. The strict inequality and $4\varepsilon<R-r$ allow a $\delta>4\varepsilon$ satisfying both inequalities in the second part. The alternative response set then preserves every acquired message. After the direction $-u$ and the safe offer at the midpoint are announced, the minimax regret exceeds $\varepsilon$. This contradicts the proposed guarantee. Hence the caps cover.

Under normalized rotation-invariant surface measure on the unit sphere, a cap of radius $\alpha$ has measure
\[
 \mu_k(\alpha)=\frac{\int_0^\alpha\sin^{k-2}\theta\,d\theta}
 {\int_0^\pi\sin^{k-2}\theta\,d\theta}.
\]
This formula follows by using the angle to the cap center as the first spherical coordinate; its surface Jacobian is proportional to $\sin^{k-2}\theta$, and the proportionality cancels in the ratio. The union bound for a cover gives $1\leq n\mu_k(\alpha)$, which is the stated exact lower bound. To see its rate without suppressing the dependence on the angular radius, write $I_k=\int_0^\pi\sin^{k-2}\theta\,d\theta$. Since $\sin\theta\leq\theta$ for nonnegative $\theta$,
\[
 n\geq(k-1)I_k\alpha^{-(k-1)}.
\]
Moreover, $\sin(\alpha/2)=\sqrt{2\varepsilon/(r+4\varepsilon)}$ and $\arcsin y\leq\pi y/2$ on $[0,1]$, so $\alpha\leq\pi\sqrt{2\varepsilon/r}$. Combining the two inequalities proves the asserted fixed-dimension lower-order rate, with the explicit positive constant $(k-1)I_k/(\pi\sqrt2)^{k-1}$.
\end{proof}

Every set used in the proof is the response set of an initial compact action menu, with zero private payoff at the center and absorbing continuation. Public randomization implements convex combinations; the queried perturbation adds the corresponding linear payoff. Thus the indistinguishable objects are admissible dynamic response economies in the stated information class. The theorem does not assert that these economies share a transition array already disclosed to the purchaser, nor that the ball construction has a finite polyhedral action representation.
